\chapter[Memory Tracing]{Memory tracing - how to compile FHI-aims using the preprocessor}
\label{appendix_memtrace}

FHI-aims features a preprocessor which can be used to add a memory tracing feature at
compile time. In its normal operation mode, it gives you a overview of the memory c
consumption at the end of the FHI-aims run, if you want also per host (node) and per process.
In its second operation mode, the tool acts as debug utility, because it gives a very
detailed account of all (de)allocations, including date, time, arraysize, current memory
consumption, peak memory consumption, the arrayname and its origin (source file and
line) for each process.\\
To use this feature, you will need in addition to the usual requirements:
\begin{itemize}
    \item a python installation (tested for Python 2.7 and 3.2)
    \item the \emph{regex} module for python (available at
        \url{https://pypi.python.org/pypi/regex}) (a user-local installation  is possible)
    \item a compiler with the \emph{SIZEOF} extension which yields the byte size of an
        object
\end{itemize}
If all these requirements are fulfilled, you can compile FHI-aims with the memory tracing
feature by following the steps listed below.
\begin{enumerate}
    \item adjust your makefile and set the variable \emph{MEMTRACE\_USE} to ``yes''
    \item set \emph{MEMTRACE\_KEEP} to ``yes'' if you want to keep the interim files
        generated by the preprocessor (suffix *.i90). The *.itypes and *.imod files will
        be kept regardless of your choice.
    \item compile FHI-aims as usual
\end{enumerate}
You now should have a version of FHI-aims that is capable of tracing the memory
consumption. The usage of this functionality is covered in the documentation of the keywords
\keyword{memtrace\_summary}, \keyword{memtrace\_use} and \keyword{memtrace\_targetdir}.